\begin{abstract}
Transformer-based architectures dominate bi-temporal \ac{BCD} from \ac{VHR} satellite imagery. However, their reliance on \ac{MHSA} modules imposes a quadratic computational and memory complexity ($\mathcal{O}(N^2)$), severely constraining scalability. This thesis presents the architectural design, implementation, and multi-temporal evaluation of \textbf{\ac{RALA-CF}}, an efficient variant of the ChangeFormer framework. The proposed methodology replaces every encoder \ac{MHSA} block with a \ac{RALA} module, mapping the computational complexity down to linear bounds ($\mathcal{O}(N)$) while preserving all auxiliary network components. On the original LEVIR-CD and DSIFN-CD datasets, \ac{RALA-CF} achieves $F_1$-scores of $89.87\%$ and $84.07\%$ respectively, preserving statistical accuracy parity within standard training variance limits of the full-rank baseline. Concurrently, our method delivers an $8\%$ reduction in attention-focused \ac{GFLOPs}, a $21.8\%$ optimization in peak \ac{VRAM} allocation, and a $50\%$ expansion in the maximum feasible mini-batch size configurations under a fixed hardware budget. To establish cross-domain generalization boundaries, this research subjects the model to an extensive six-dataset benchmark, adding LEVIR-CD+, SYSU-CD, WHU-CD, and S2Looking. The quantitative evaluations reveal a performance gradient governed by acquisition geometry and scene heterogeneity. The model transfers reliably to nadir-view aerial datasets with high-precision annotations, such as WHU-CD ($F_1 = 87.85\%$), but experiences structural degradation under oblique side-looking perspectives, exemplified by the S2Looking benchmark ($F_1 = 60.39\%$). These findings frame the proposed architecture as a highly practical efficiency contribution that democratizes large-scale remote sensing inference on standard computing infrastructure.

\vspace{0.8cm}
\noindent{\textbf{Keywords:}} Building Change Detection, Linear Attention, Remote Sensing, Computational Efficiency, Satellite Imagery, Vision Transformers, Benchmarking.
\end{abstract}
